\begin{lemma}
Let $\mathcal H$ be a $3$-uniform hypergraph of maximum degree at most $D$,
let $e\in E(\mathcal H)$ intersect exactly $3(D-1)$ other edges, and let
$G$ be the subgraph of the line graph $L(\mathcal H)$ induced on
$N_{L(\mathcal H)}(e)$.  Put
\[
X=X(e,\mathcal H)=
\left(\bigcup\{h\in E(\mathcal H):h\cap e\ne\emptyset\}\right)\setminus e
\]
and
\[
w_{\mathcal H,e}(X)=
\sum_{x\in X}\sum_{v<v'\in e}
\deg_{\mathcal H}(v,x)\deg_{\mathcal H}(v',x).
\]
After fixing an arbitrary order on the three vertices of $e$, define
\[
Y_{\mathcal H}(e)=
\sum_{v<v'\in e}
\sum_{\substack{e_1\ne e,\ v\in e_1\\ e_2\ne e,\ v'\in e_2}}
\max\{0, |e_1\cap e_2|-1\}.
\]
Then $Y_{\mathcal H}(e)\ge 0$ and
\[
|E(G)|=3\binom{D-1}{2}+w_{\mathcal H,e}(X)-Y_{\mathcal H}(e).
\]
\end{lemma}
\begin{proof}
The incidence equality argument below is formalized, without any simplicity
assumption, by \noderef{SaturatedNeighborClassPartition}.
By \noderef{UniformHypergraph}, the edge $e$ has exactly three vertices; fix
an ordering $e=\{v_1,v_2,v_3\}$.  By \noderef{HypergraphDegree} and
\noderef{MaxDegreeAtMost}, the assertion that $\mathcal H$ has maximum
degree at most $D$ means that each $v_i$ is contained in at most $D$ edge
indices of $\mathcal H$, hence in at most $D-1$ edge indices other than
$e$.  The hypothesis says that the number of edges other than $e$ meeting
$e$ is exactly $3(D-1)$.

We first record the consequence that every edge $h\ne e$ meets $e$ in at
most one vertex.  If some $h\ne e$ contained two vertices of $e$, then the
sum over $i=1,2,3$ of the numbers of edges other than $e$ containing $v_i$
would still be at most $3(D-1)$, but $h$ would be counted at least twice.
Therefore the number of distinct edges meeting $e$ would be at most
$3(D-1)-1$, contradicting the assumed value $3(D-1)$.  Thus no neighboring
edge belongs to two of the three classes
\[
C_i=\{h\in E(\mathcal H):h\ne e,\ v_i\in h\}.
\]
Since the union of the disjoint classes $C_1,C_2,C_3$ has size $3(D-1)$
and each class has size at most $D-1$, each $C_i$ has exactly $D-1$ edges.

By \noderef{LineGraphOfHypergraph}, two vertices of the line graph
$L(\mathcal H)$ are adjacent exactly when the corresponding hyperedges are
distinct and intersect.  Therefore the induced graph $G$ on
$N_{L(\mathcal H)}(e)$ has vertex set $C_1\cup C_2\cup C_3$.  Within each
$C_i$, every two distinct hyperedges both contain $v_i$, so they are
adjacent in $L(\mathcal H)$ and hence in $G$.  These three disjoint cliques
contribute
\[
3\binom{D-1}{2}
\]
edges to $G$.

It remains to count the edges of $G$ between different classes.  Fix
$i<j$.  For $e_1\in C_i$ and $e_2\in C_j$, the common vertices
$e_1\cap e_2$ cannot contain $v_i$ or $v_j$, because $e_2$ does not contain
$v_i$ and $e_1$ does not contain $v_j$ by the preceding paragraph.  Hence
$e_1$ and $e_2$ are adjacent in $G$ exactly when they share at least one
vertex of $X$.

In the formal expression for $w$, the sum is over two-element subsets of
$e$. The bijection sending an increasing pair $(v,v')$ to $\{v,v'\}$,
and the associated product-sum identity, are supplied by
\noderef{FiniteUnorderedPairProductSum}. Thus this expression is exactly
the sum over the pairs of classes used here.

For a fixed vertex $x\in X$, the number of ordered choices
$(e_1,e_2)\in C_i\times C_j$ with $x\in e_1\cap e_2$ is
$\deg_{\mathcal H}(v_i,x)\deg_{\mathcal H}(v_j,x)$, where
$\deg_{\mathcal H}(v,x)$ denotes the number of edge indices of
$\mathcal H$ containing both $v$ and $x$.  Indeed, the first factor chooses
the edge in $C_i$ containing $x$, and the second factor chooses the edge in
$C_j$ containing $x$.  Summing over all $x\in X$ therefore counts a fixed
pair $(e_1,e_2)\in C_i\times C_j$ once for each common outside vertex
$x\in e_1\cap e_2\cap X$.  As observed above, common vertices of such a
cross-class pair cannot be $v_i$ or $v_j$, and the only vertices of $e$ in
$e_1\cup e_2$ are $v_i$ and $v_j$; hence every common vertex of
$e_1$ and $e_2$ lies in $X$.  Thus the sum over $x\in X$ counts the pair
exactly $|e_1\cap e_2|$ times.

The true line-graph edge should be counted once when this intersection is
nonempty and zero times otherwise.  Therefore one subtracts, for this
pair, precisely the excess
\[
|e_1\cap e_2|-\mathbf 1[e_1\cap e_2\ne\emptyset]
=\max\{0,|e_1\cap e_2|-1\}.
\]
Consequently the number of cross-edges between $C_i$ and $C_j$ is
\[
\sum_{x\in X}\deg_{\mathcal H}(v_i,x)\deg_{\mathcal H}(v_j,x)
-
\sum_{\substack{e_1\in C_i\\ e_2\in C_j}}
\max\{0,|e_1\cap e_2|-1\}.
\]
This double count and the integer excess correction are also stated together
in \noderef{FiniteBipartiteIntersectionEdgeCount}, applied with the indexed
edge sets $C_i,C_j$ and the outside vertex set $X$.

Summing this identity over the three pairs $i<j$ gives the total number of
cross-edges as $w_{\mathcal H,e}(X)-Y_{\mathcal H}(e)$.  Adding the
within-clique contribution gives
\[
|E(G)|=3\binom{D-1}{2}+w_{\mathcal H,e}(X)-Y_{\mathcal H}(e).
\]
Finally, every summand in the definition of $Y_{\mathcal H}(e)$ is a maximum
of $0$ and another integer, so each summand is nonnegative and therefore
$Y_{\mathcal H}(e)\ge0$.
\end{proof}
